\documentclass[11pt]{amsart}
\usepackage{amsmath,amssymb,amsthm,mathtools}
\usepackage[margin=1.05in]{geometry}
\usepackage{booktabs}

\numberwithin{equation}{section}
\newtheorem{theorem}{Theorem}[section]
\newtheorem{proposition}[theorem]{Proposition}
\newtheorem{lemma}[theorem]{Lemma}
\newtheorem{corollary}[theorem]{Corollary}
\theoremstyle{definition}
\newtheorem{definition}[theorem]{Definition}
\newtheorem{hypothesis}[theorem]{Hypothesis}
\theoremstyle{remark}
\newtheorem{remark}[theorem]{Remark}

\newcommand{\Z}{\mathbb{Z}}
\newcommand{\Q}{\mathbb{Q}}
\newcommand{\R}{\mathbb{R}}
\newcommand{\N}{\mathbb{N}}
\newcommand{\C}{\mathbb{C}}
\newcommand{\HH}{\mathbb{H}}
\newcommand{\cH}{\mathcal{H}}
\newcommand{\cM}{\mathcal{M}}
\newcommand{\cT}{\mathcal{T}}
\newcommand{\cR}{\mathcal{R}}
\newcommand{\cG}{\mathcal{G}}
\newcommand{\Nrd}{\operatorname{Nrd}}
\newcommand{\Ad}{\operatorname{Ad}}
\newcommand{\id}{\mathrm{id}}
\newcommand{\A}{\mathbb{A}}
\newcommand{\op}{\mathrm{op}}

\PassOptionsToPackage{hyphens}{url}
\usepackage[hidelinks,bookmarks=false]{hyperref}
\hypersetup{pdftitle={The Ramanujan bound is the tempered bound: a no-go theorem for condition
(C4)}}

\title[The Ramanujan bound is the tempered bound]{The Ramanujan bound is the tempered bound:\\
the Hecke operator on the Bruhat--Tits tree, and\\ a no-go theorem for condition (C4)}

\author{Ruqing Chen}
\address{GUT Geoservice Inc., Montreal, Canada}
\email{ruqing@hotmail.com}
\date{}

\raggedbottom
\begin{document}

\begin{abstract}
Condition (C4) of Paper X of this series asks whether the measured groupoid of the Hurwitz
quaternion Bost--Connes system is non-amenable, equivalently whether its factor escapes the
injectivity that erased the arithmetic in every earlier model. Paper XI reduced the question to a
spectral property of a transfer operator and found the naive operator unusable, the range
projections being nested rather than orthogonal; Paper XII reduced it further to a single
representation of the norm-one group and warned that a spectral gap would not suffice. We settle
the question.

We first build the operator the programme wanted: the Hecke operator
$T_p=\sum_{k=1}^{p+1}\pi(a_k)$, the sum of the Koopman unitaries of the $p+1$ Hurwitz quaternions
of reduced norm $p$, which is the adjacency operator of the $(p+1)$-regular Bruhat--Tits tree of
$\mathrm{GL}_2(\Q_p)$ transported to the Maharam extension. It involves no projections and so has
none of the pathology of the transfer operator.

We then compute its spectrum exactly. For every finite set $F$ of primes the subgroup
$G_F=\HH(\Z[1/F])^\times$ is discrete in $\prod_{p\in F}\HH(\Q_p)^\times$ and acts \emph{properly
and freely} on a conull set, because every element of $\HH(\Q)^\times$ is a unit at almost every
place. Hence the Koopman representation restricted to $G_F$ is a multiple of the regular
representation, and, with $\varepsilon_p$ the central unitary of $p$ (so that
$\varepsilon_p^{-1/2}T_p$ is the self-adjoint Hecke operator),
\[
\mathrm{sp}\bigl(\varepsilon_p^{-1/2}T_p\restriction L^2_0\bigr)=[-2\sqrt p,\,2\sqrt p],\qquad \|T_p\|=2\sqrt p .
\]
The Ramanujan bound is attained, not merely respected. It is also worthless for the purpose:
$2\sqrt p$ is the norm of the adjacency operator in the \emph{regular} representation of a free
group, so the gap between $2\sqrt p$ and the trivial eigenvalue $p+1$ is Kesten's gap and measures
the non-amenability of the \emph{group}, not of the \emph{action}. By Kuhn's theorem an amenable
action has tempered Koopman representation; attaining the Ramanujan bound proves temperedness, and
temperedness is the signature of amenability, not an obstruction to it.

The same properness gives the verdict directly, with no spectral theory at all. The orbit relation
is the increasing union over finite $F$ of the smooth relations of the proper actions of $G_F$,
hence hyperfinite; the groupoid is amenable and the factor is injective. \textbf{(C4) is false},
the Hurwitz factor collapses to $R_\infty$, $R_\lambda$ or a Krieger factor like every abelian
system before it, and by Paper IX it retains nothing of the arithmetic. The argument uses only
that each element of the monoid is a unit at almost every place, so it closes the search of
Paper X in full: conditions (C1)--(C3) are incompatible with (C4) for every order in every
finite-dimensional $\Q$-algebra.
\end{abstract}

\maketitle

\section{Setting and statements}\label{sec:setup}

$\HH$ is the rational Hamilton quaternion algebra, ramified at $2$ and $\infty$;
$\cH\subset\HH(\Q)$ is the Hurwitz order, $P=\cH\setminus\{0\}$ the multiplicative Ore monoid of
Paper X, $\Nrd$ the reduced norm and $G=PP^{-1}=\HH(\Q)^\times$ the group of right fractions
\cite{Vigneras}. For a prime $p$ write $\cH_p=\cH\otimes\Z_p$; for odd $p$ one has
$\cH_p\cong M_2(\Z_p)$ and the right ideals of $\cH_p$ of reduced norm $p$ are the $p+1$ neighbours
of the base vertex in the Bruhat--Tits tree $\cT_p$ of $\mathrm{GL}_2(\Q_p)$ \cite{Serre}. The
localized Bost--Connes framework in which the system sits is that of \cite{Main}; the ingredients
we use from Papers IX--XII of the series are re-proved here.

\begin{hypothesis}[The model]\label{hyp:model}
$\Omega=\widehat\cH=\prod_p\cH_p$ carries a $\mathrm{KMS}_\beta$ measure $\nu_\beta$ of the
Nica--Toeplitz algebra $\cT(P)$, $\widetilde\Omega=\varinjlim(\Omega,\cdot a)=\HH(\A_f)$ is the
dilation, $G$ acts on $\widetilde\Omega$ by right translation, and $\widetilde\nu_\beta$ is the
extension of $\nu_\beta$, so that
$\widetilde\nu_\beta(Ea)=\Nrd(a)^{-\beta}\widetilde\nu_\beta(E)$. The Maharam extension is
\begin{equation}\label{eq:maharam}
G\curvearrowright\bigl(\widetilde\Omega\times\R,\ \widetilde\nu_\beta\otimes e^{s}ds\bigr),
\qquad a\cdot(x,s)=\bigl(xa,\ s+\beta\log\Nrd(a)\bigr),
\end{equation}
which preserves the displayed infinite measure. Condition (C4) of Paper X asks whether the
measured groupoid $G\ltimes(\widetilde\Omega,\widetilde\nu_\beta)$ is non-amenable, equivalently
--- by Takesaki duality, as in Paper XI --- whether \eqref{eq:maharam} is non-amenable,
equivalently whether $M_\beta=\pi_\varphi(\cT(P))''$ fails to be injective.
\end{hypothesis}

Nothing below uses the value of $\beta$, the existence of a $\mathrm{KMS}$ state in any particular
temperature range, or the type of $M_\beta$. Only the following is used, and it is a property of
the arithmetic and not of the dynamics.

\begin{lemma}[Finite support]\label{lem:support}
Every $a\in G=\HH(\Q)^\times$ lies in $\cH_q^\times$ for all but finitely many primes $q$.
\end{lemma}

\begin{proof}
Write $a=b/m$ with $b\in\cH$ and $m\in\Z_{>0}$. For $q\nmid m\Nrd(b)$ one has $b\in\cH_q$ and
$\Nrd(b)\in\Z_q^\times$, so $b\bar b=\Nrd(b)$ is a unit and $b\in\cH_q^\times$; and
$m\in\Z_q^\times$. Only the finitely many $q$ dividing $m\Nrd(b)$ are excluded.
\end{proof}

For a finite set $F$ of primes put
\[
G_F=\bigl\{a\in G:\ a\in\cH_q^\times\ \text{for every } q\notin F\bigr\}=\HH\bigl(\Z[1/F]\bigr)^\times,
\qquad L_F=\prod_{p\in F}\HH(\Q_p)^\times .
\]
By Lemma~\ref{lem:support}, $G=\bigcup_FG_F$, the union being increasing and countable.

\medskip
The three results are the following.

\begin{theorem}[Local properness]\label{thm:proper}
For every finite $F$, $G_F$ is a discrete subgroup of $L_F$ and its action on $\widetilde\Omega$
is free and proper on a conull invariant set. Consequently the orbit relation $\cR_F$ is smooth,
and the Koopman representation of the Maharam extension restricted to $G_F$ is unitarily
equivalent to a multiple of the regular representation $\lambda_{G_F}$.
\end{theorem}

\begin{theorem}[The spectrum of the Hecke operator]\label{thm:spec}
Let $p$ be an odd prime, $a_1,\dots,a_{p+1}\in\cH$ generators of the $p+1$ right ideals of reduced
norm $p$, and $T_p=\sum_{k=1}^{p+1}\pi(a_k)$ the associated Hecke operator on
$L^2\bigl(\widetilde\Omega\times\R\bigr)$ of Definition~\ref{def:hecke}. Then $T_p$ is the
adjacency operator of $\cT_p$ in the regular representation up to the central unitary
$\varepsilon_p=\pi(p)$, and on the orthogonal complement
$L^2_0$ of the invariant vectors the self-adjoint operator $\varepsilon_p^{-1/2}T_p$ of
Lemma~\ref{lem:hecke}(1) satisfies
\[
\mathrm{sp}\bigl(\varepsilon_p^{-1/2}T_p\restriction L^2_0\bigr)=\bigl[-2\sqrt p,\ 2\sqrt p\,\bigr],
\qquad \bigl\|T_p\restriction L^2_0\bigr\|=2\sqrt p .
\]
($T_p$ itself is normal but not self-adjoint, $T_p^*=\varepsilon_p^{-1}T_p$; its spectrum is
the disc swept by the central phases, and the invariant statement is the one displayed.)
The Ramanujan bound is attained, and the trivial eigenvalue $p+1$ is absent from $L^2_0$.
\end{theorem}

\begin{theorem}[No-go for (C4)]\label{thm:nogo}
The orbit relation $\cR_G=\bigcup_F\cR_F$ is hyperfinite. Hence the measured groupoid of
Hypothesis~\ref{hyp:model} is amenable, (C4) is false, and $M_\beta$ is injective: it is
$R_\infty$, a Powers factor $R_\lambda$, a Krieger factor, or type I. By Paper IX it therefore
retains no arithmetic invariant of the system.
\end{theorem}

Theorem~\ref{thm:nogo} does not use Theorem~\ref{thm:spec}; the spectral computation is included
because it identifies precisely what the programme of Paper XI was measuring, and why the
measurement could not have decided the question. Section~\ref{sec:general} extends the no-go to
every candidate monoid admissible under (C1)--(C3).

\section{The Hecke operator}\label{sec:hecke}

Paper XI worked with the transfer operator $\cM_p(x)=\frac1{p+1}\sum_kv_{a_k}^*xv_{a_k}$ built
from the isometries of $\cT(P)$. Its defect is structural: the range projections
$e_{a_k}=v_{a_k}v_{a_k}^*$ are nested, $e_{a_j}e_{a_k}=e_p\ne0$, so $\cM_p$ is not a Markov
operator of the tree and forcing orthogonality collapses the algebra. The remedy is to work on the
Maharam extension, where the generators act by \emph{unitaries} and no projection occurs.

\begin{definition}\label{def:hecke}
Let $\pi$ be the Koopman representation of \eqref{eq:maharam},
\[
\bigl(\pi(a)f\bigr)(x,s)=f\bigl(xa,\ s+\beta\log\Nrd(a)\bigr),
\qquad a\in G,
\]
a unitary representation of $G$ on $L^2(\widetilde\Omega\times\R,\widetilde\nu_\beta\otimes e^sds)$
because the Maharam extension preserves the measure. For odd $p$ fix representatives
$a_1,\dots,a_{p+1}\in\cH$ of the $p+1$ right ideals of reduced norm $p$ and set
\[
T_p=\sum_{k=1}^{p+1}\pi(a_k).
\]
\end{definition}

\begin{lemma}\label{lem:hecke}
With $\varepsilon_p=\pi(p)$ the unitary of the central element $p$:
\begin{enumerate}
\item the set $\{a_k\}$ may be chosen closed under $a\mapsto\bar a$ modulo right units, and then
$T_p^*=\varepsilon_p^{-1}T_p$, so $\varepsilon_p^{-1/2}T_p$ is self-adjoint for either square root;
\item $T_p^2=T_{p^2}+(p+1)\varepsilon_p$, where $T_{p^2}$ is the sum over the $p(p+1)$ right
ideals of reduced norm $p^2$ not containing $p\cH$;
\item under the identification of Theorem~\ref{thm:proper}, $T_p$ is the adjacency operator of the
$(p+1)$-regular tree $\cT_p$ acting on $\ell^2$ of its vertex set.
\end{enumerate}
\end{lemma}

\begin{proof}
(1) $a_k\bar a_k=p$, so $\pi(\bar a_k)=\varepsilon_p\pi(a_k)^{-1}=\varepsilon_p\pi(a_k)^*$;
summing over a self-conjugate set of representatives gives the identity. (2) Expanding
$T_p^2=\sum_{j,k}\pi(a_ja_k)$, the products $a_ja_k$ of reduced norm $p^2$ either generate an ideal
not containing $p\cH$, and each such ideal arises exactly once, or satisfy $a_ja_k\in p\cH$, which
happens for exactly one $k$ for each $j$: $a_ja_k\in p\cH$ forces $a_k\cH=\bar a_j\cH$, and
with the self-conjugate choice of (1) then $a_k=\bar a_j$ exactly, so the product is
$a_j\bar a_j=p$ and contributes $\varepsilon_p$ exactly; there are $p+1$ such pairs. (3) The vertices of $\cT_p$ at distance $n$ from the base vertex are
the right ideals of $\cH_p$ of reduced norm $p^n$ that are not divisible by $p$ beyond the
necessary amount, and $a\mapsto$ (the ideal it generates) intertwines right multiplication by the
$a_k$ with the adjacency of $\cT_p$; the identification of the Hilbert space is
Corollary~\ref{cor:regular}. (Vertices are ideal classes up to homothety, i.e.\ after
trivializing the central $\langle p\rangle$; on the Koopman space itself the central unitary
$\varepsilon_p$ survives, which is why the self-adjoint operator is
$\varepsilon_p^{-1/2}T_p$.)
\end{proof}

\begin{remark}
By Lubotzky--Phillips--Sarnak the group generated by $a_1,\dots,a_{p+1}$ modulo the centre and the
$24$ units is free of rank $(p+1)/2$ and acts simply transitively on the vertices of $\cT_p$; the
finite quotients by congruence subgroups are the Ramanujan graphs $X^{p,q}$. This is the algebraic
rigidity the present programme hoped to exploit.
\end{remark}

\section{Local properness and the regular representation}\label{sec:proper}

\begin{proof}[Proof of Theorem~\ref{thm:proper}]
\emph{Discreteness.} Let $a_n\in G_F$ with $a_n\to1$ in $L_F$. Then $\Nrd(a_n)\to1$ in
$\prod_{p\in F}\Q_p^\times$, and $\Nrd(a_n)\in\Z[1/F]^\times=\{\pm\prod_{p\in F}p^{e_p}\}$, so for
large $n$ all exponents vanish and $\Nrd(a_n)=1$. Thus $a_n$ lies in the norm-one group
$\HH(\Z[1/F])^1$, which is a lattice in $\prod_{p\in F}\HH(\Q_p)^1$ because $\HH$ is definite and
$\HH(\R)^1=SU(2)$ is compact; in particular it is discrete, so $a_n=1$ for large $n$.

\emph{Properness and freeness.} Write $\widetilde\Omega=X_F\times Y_F$ with
$X_F=\prod_{p\in F}\HH(\Q_p)$ and $Y_F$ the remaining coordinates. The locus where the component at $p$
has zero reduced norm is $\widetilde\nu_\beta$-null: an element of $\cH_p$ of zero reduced
norm lies, for every $n$, in one of the right ideals of reduced norm $p^n$, of which there
are $\sigma(p^n)=1+p+\cdots+p^n$, each carried by the scaling relation to measure
$p^{-n\beta}$ times that of the base cylinder (an element of reduced norm $p^n$ is a unit at
every other place, by Lemma~\ref{lem:support}, so the cylinder is an honest translate); the
bound $\sigma(p^n)p^{-n\beta}\to0$ holds for $\beta>1$, in particular throughout the Gibbs
range $\beta>2$ of Paper X. So $X_F^\times=\prod_{p\in F}\HH(\Q_p)^\times$ is conull and invariant.
By Lemma~\ref{lem:support} and the definition of $G_F$, an element $a\in G_F$ acts on $Y_F$ through
the compact group $\prod_{q\notin F}\cH_q^\times$ and on $X_F^\times$ by right translation by the
discrete subgroup $G_F\le L_F$. Right translation of a locally compact group by a discrete subgroup
is free and proper, and properness in one factor of a product is properness; so the action on
$X_F^\times\times Y_F$ is free and proper.

\emph{Smoothness and the Koopman representation.} A free proper action of a countable group admits
a measurable fundamental domain $D$, so $\widetilde\Omega\times\R=\bigsqcup_{a\in G_F}(D\cdot a)$
up to null sets and $\cR_F$ is smooth. Trivializing the Radon--Nikodym derivative by half-densities
gives a unitary
$L^2(\widetilde\Omega\times\R)\cong L^2(D)\otimes\ell^2(G_F)$ intertwining $\pi\restriction G_F$
with $1\otimes\lambda_{G_F}$.
\end{proof}

\begin{corollary}\label{cor:regular}
$\pi\restriction G_F\cong\dim L^2(D)\cdot\lambda_{G_F}$ for every finite $F$. In particular
$\pi\restriction G_F$ is tempered, and for every $\xi$ in the group algebra
$\|\pi(\xi)\|=\|\lambda_{G_F}(\xi)\|$.
\end{corollary}

\section{The spectrum, and what it does not prove}\label{sec:spec}

\begin{proof}[Proof of Theorem~\ref{thm:spec}]
Take $F=\{p\}$, so $T_p=\pi(\xi)$ with $\xi=\sum_ka_k$ in the group algebra of $G_{\{p\}}$. By
Corollary~\ref{cor:regular}, $\|T_p\|=\|\lambda(\xi)\|$ and more precisely
$\mathrm{sp}(T_p)=\mathrm{sp}_{L(G_{\{p\}})}(\xi)$, the spectrum in the group von Neumann algebra.
Modulo the centre and the $24$ units, the $a_k$ generate a group $\Gamma$
acting simply transitively on the vertices of $\cT_p$ --- free of rank $(p+1)/2$ when
$p\equiv1\ (\mathrm{mod}\ 4)$ \cite{LPS}, a free product of cyclic groups in general ---
and $\xi$ maps to the sum of the
$p+1$ generating translations, whose image under $\lambda_\Gamma$ is the adjacency operator of the
Cayley graph of $\Gamma$, which is $\cT_p$ \cite{Serre}. To pass from $\ell^2(G_{\{p\}})$ to
$\ell^2$ of the vertex set, decompose over the characters of the central subgroup generated
by $p$ and the finite unit group: in each fibre $T_p$ acts as a phase-twisted tree adjacency
operator, and on a tree every phase twist is gauge-trivial, so the fibre operator is
unitarily a phase times the tree adjacency, and the self-adjoint $\varepsilon_p^{-1/2}T_p$
has in every fibre the spectrum of the tree adjacency itself. By Kesten's theorem \cite{Kesten} the
spectrum of the adjacency
operator of the $(p+1)$-regular tree on $\ell^2$ of its vertices is exactly
$[-2\sqrt p,2\sqrt p\,]$, with norm $2\sqrt p$ and no eigenvalues. Invariant vectors would carry
the eigenvalue $p+1$ and do not lie in $\ell^2$, so they are absent from $L^2_0$.
\end{proof}

\begin{corollary}[There is a gap, and it is Kesten's]\label{cor:gap}
$\|T_p\restriction L^2_0\|=2\sqrt p<p+1$, a gap of relative size $1-2\sqrt p/(p+1)$:
\[
\begin{array}{lrrrrrr}
p & 3 & 5 & 7 & 11 & 13 & 101\\
2\sqrt p & 3.464 & 4.472 & 5.292 & 6.633 & 7.211 & 20.10\\
p+1 & 4 & 6 & 8 & 12 & 14 & 102\\
\text{gap} & 0.134 & 0.255 & 0.339 & 0.447 & 0.485 & 0.803
\end{array}
\]
This gap is present in the regular representation of $\Gamma$ itself, and its existence is
equivalent, by Kesten's criterion \cite{Kesten}, to the non-amenability of the free group
$\Gamma$. It is
therefore a statement about the \emph{group} $G$, not about the \emph{action} of $G$ on
$\widetilde\Omega$.
\end{corollary}

\begin{theorem}[The proposed equivalence is vacuous]\label{thm:vacuous}
The following are equivalent, and none of them implies non-amenability of the action:
\begin{enumerate}
\item $\|T_p\restriction L^2_0\|=2\sqrt p$, the Ramanujan bound, for every odd $p$;
\item $\pi\restriction G_F$ is weakly contained in $\lambda_{G_F}$ for every finite $F$, i.e.\ the
Koopman representation is tempered;
\item the free group $\Gamma$ is non-amenable.
\end{enumerate}
Indeed (1) and (2) \emph{hold}, by Theorem~\ref{thm:spec} and Corollary~\ref{cor:regular}, while
the action is amenable by Theorem~\ref{thm:nogo}.
\end{theorem}

\begin{proof}
(1)$\Leftrightarrow$(3) is Kesten's criterion \cite{Kesten}. (2)$\Rightarrow$(1) because temperedness gives
$\|\pi(\xi)\|\le\|\lambda(\xi)\|$ and the reverse inequality is Corollary~\ref{cor:regular};
(1)$\Rightarrow$(2) for the algebra generated by the $a_k$, and the general case is
Corollary~\ref{cor:regular} again. The last assertion is the content of the paper.
\end{proof}

\begin{remark}[The orientation of the argument]\label{rem:orientation}
Kuhn's theorem \cite{Kuhn}, quoted in Paper XII, says that an amenable action has tempered
Koopman representation. Establishing the Ramanujan bound therefore establishes temperedness, which is a
\emph{necessary condition for amenability}, not an obstruction to it. To prove (C4) by spectral
means one would have to exhibit spectrum of $T_p$ \emph{outside} $[-2\sqrt p,2\sqrt p]$ on
$L^2_0$ --- that is, a \emph{violation} of the Ramanujan property, the exact opposite of what the
arithmetic of the Hurwitz order provides. The bound $2\sqrt p$ is not a strong constraint that the
programme could exploit; it is the Plancherel measure of $\mathrm{PGL}_2(\Q_p)$ announcing that
there is nothing beyond the regular representation to find.
\end{remark}

\section{The no-go theorem}\label{sec:nogo}

\begin{proof}[Proof of Theorem~\ref{thm:nogo}]
By Lemma~\ref{lem:support} the sets $\cR_F$, $F$ finite, form an increasing sequence of
subrelations of $\cR_G$ with union $\cR_G$. By Theorem~\ref{thm:proper} each $\cR_F$ is smooth,
hence hyperfinite. An increasing union of hyperfinite equivalence relations is hyperfinite
\cite{Dye,CFW}, so $\cR_G$ is hyperfinite, hence amenable \cite{Zimmer}. Amenability of the orbit
relation is amenability of the measured groupoid, since the isotropy is trivial on the conull set
of Theorem~\ref{thm:proper} and the relation is the one attached to the groupoid by \cite{FM};
amenability of the groupoid is equivalent to injectivity of the associated von Neumann algebra by
\cite{CFW,ADR}. The classification of injective factors then leaves only
$R_\infty$, $R_\lambda$, a Krieger factor and the type I case, and the reconstruction failure is
Paper IX.
\end{proof}

\begin{remark}[Why no spectral input was needed]\label{rem:whyno}
The proof does not see the tree, the Ramanujan property or the non-amenability of $G$. It uses one
arithmetic fact, Lemma~\ref{lem:support}: an element of $G$ has a reduced norm with finitely many
prime factors, so it acts at finitely many places at a time. The group never acts \emph{all at
once}. Non-amenability of $G$ is a statement about the group acting on itself; the groupoid only
ever sees the finite-level subgroups $G_F$, each of which acts properly.
\end{remark}

\begin{remark}[The $F_2$ example was the mechanism]\label{rem:F2}
Paper XII observed, as a warning about Zimmer's theorem, that the translation action of $F_2$ on
itself preserves counting measure, is essentially free, and is nonetheless amenable, being proper.
Theorem~\ref{thm:proper} shows that this is not an analogy: the Hurwitz action \emph{is} an
increasing union of translation actions of discrete groups on locally compact groups. The
counterexample and the system coincide.
\end{remark}

\section{The search of Paper X is empty}\label{sec:general}

\begin{theorem}[General no-go]\label{thm:general}
Let $A$ be a finite-dimensional $\Q$-algebra with reduced norm $\Nrd$, $\Lambda\subset A$ an
order, $P\subseteq\Lambda\setminus\{0\}$ a right Ore, right LCM monoid with $G=PP^{-1}\le A^\times$
and $N=|\Nrd|$ the norm of Paper X. Suppose $G$ acts by right translation on the finite-adelic
space $\widetilde\Omega=A(\A_f)$ with a quasi-invariant measure scaling by $N^{-\beta}$, and that
$G$ is discrete in $\prod_{p\in F}A(\Q_p)^\times$ for every finite $F$ --- automatic when
the orders of $A$ have finite unit groups, as for definite quaternion algebras over $\Q$;
for orders with infinite unit groups it can fail, already for a real quadratic field, where
the closure of the unit group in the $p$-adic factors is the Leopoldt object of Papers
XXII--XXIII, and the hypothesis must then be checked. Then the
orbit relation is hyperfinite and the associated factor is injective, for every $\beta$.
\end{theorem}

\begin{proof}
Lemma~\ref{lem:support} holds verbatim: $a\in G$ has $\Nrd(a)$ a rational number, so $a$ is a unit
of $\Lambda\otimes\Z_q$ for all but the finitely many $q$ dividing its numerator and denominator.
The proof of Theorem~\ref{thm:proper} then applies with $\HH$ replaced by $A$, and the proof of
Theorem~\ref{thm:nogo} follows.
\end{proof}

\begin{corollary}\label{cor:closed}
Conditions (C1)--(C3) of Paper X are incompatible with (C4). Any monoid satisfying (C1) sits inside
an order in a finite-dimensional $\Q$-algebra, by (C2) and the finiteness of the norm; (C3)
forces the unit group $P^\times$ to be finite (Paper X, its Remark 1.5: an infinite unit
group makes $\zeta_P$ diverge identically), so the norm-one group of the order is finite,
the archimedean norm-one group is compact, and the discreteness hypothesis of
Theorem~\ref{thm:general} holds; and every
element of the monoid is a unit at almost every place; so the resulting groupoid is amenable and the factor
is injective. The search for a non-injective arithmetic Bost--Connes factor cannot succeed by
enlarging the index monoid.
\end{corollary}

\begin{remark}[What would be required]\label{rem:escape}
Theorem~\ref{thm:general} isolates the obstruction exactly: the finite support of the norm. An
index monoid escaping it would have to contain elements that are non-units at infinitely many
places, hence of infinite norm, contradicting the multiplicativity of $N$ into $[2,\infty)$
required by (C1). Within the framework of Paper X, therefore, no escape exists. Enlargements
outside it --- an index set that is not of arithmetic type, or a coefficient algebra that is not a
function algebra --- are excluded by Paper II and Paper XXVII respectively.
\end{remark}

\section{Verdict}\label{sec:verdict}

\[
\begin{array}{@{}l@{\qquad}l@{\qquad}l@{}}
\text{Statement} & \text{Status} & \text{Where}\\\hline
T_p\ \text{is the tree adjacency operator} & \text{true} & \text{Lem.~\ref{lem:hecke}}\\
G_F\ \text{acts properly and freely} & \text{true} & \text{Thm.~\ref{thm:proper}}\\
\pi\restriction G_F\cong\infty\cdot\lambda_{G_F} & \text{true} & \text{Cor.~\ref{cor:regular}}\\
\mathrm{sp}(\varepsilon_p^{-1/2}T_p\restriction L^2_0)=[-2\sqrt p,2\sqrt p] & \text{true, attained} & \text{Thm.~\ref{thm:spec}}\\
\text{a spectral gap exists} & \text{true, and it is Kesten's} & \text{Cor.~\ref{cor:gap}}\\
\text{the gap implies (C4)} & \textbf{false} & \text{Thm.~\ref{thm:vacuous}}\\
\text{the orbit relation is hyperfinite} & \text{true} & \text{Thm.~\ref{thm:nogo}}\\
\text{(C4)} & \textbf{false} & \text{Thm.~\ref{thm:nogo}}\\
\text{some other monoid satisfies (C4)} & \textbf{false} & \text{Thm.~\ref{thm:general}}
\end{array}
\]

The Hurwitz system has everything the programme asked for. Its group of fractions is
non-amenable; it carries free subgroups of every finite rank; it generates the Bruhat--Tits trees
whose congruence quotients are the Ramanujan graphs; the Hecke operator on the Maharam extension
is exactly the tree adjacency operator, with no nested projections and no Cuntz obstruction; and
its spectrum attains the Ramanujan bound. None of it matters. The bound $2\sqrt p$ is the norm of
the same operator in the regular representation, so attaining it certifies that the Koopman
representation is regular --- and a regular Koopman representation is what a free proper action
has. The rigidity is real and it belongs to the group; the action never inherits it, because an
element of an arithmetic monoid is a unit at almost every place and the group acts only finitely
many places at a time.

The Bost--Connes phase transition therefore survives every enlargement of the index monoid
compatible with its own axioms, and the factor does not. The arithmetic lives in the
$C^*$-algebra with its flow, in the $\mathrm{KMS}$ simplex, and in the obstruction group; the von
Neumann algebra is, and remains, hyperfinite.

\begin{thebibliography}{99}
\bibitem{ADR} C. Anantharaman-Delaroche, J. Renault, \emph{Amenable Groupoids}, Monogr.
Enseign. Math. 36, Geneva, 2000.
\bibitem{CFW} A. Connes, J. Feldman, B. Weiss, \emph{An amenable equivalence relation is generated
by a single transformation}, Ergodic Theory Dynam. Systems 1 (1981), 431--450.
\bibitem{Dye} H. A. Dye, \emph{On groups of measure preserving transformations I}, Amer. J. Math.
81 (1959), 119--159.
\bibitem{FM} J. Feldman, C. C. Moore, \emph{Ergodic equivalence relations, cohomology, and von
Neumann algebras I, II}, Trans. Amer. Math. Soc. 234 (1977), 289--324, 325--359.
\bibitem{Kesten} H. Kesten, \emph{Symmetric random walks on groups}, Trans. Amer. Math. Soc. 92
(1959), 336--354.
\bibitem{Kuhn} M. G. Kuhn, \emph{Amenable actions and weak containment of some representations of
discrete groups}, Proc. Amer. Math. Soc. 122 (1994), 751--757.
\bibitem{LPS} A. Lubotzky, R. Phillips, P. Sarnak, \emph{Ramanujan graphs}, Combinatorica 8
(1988), 261--277.
\bibitem{Serre} J.-P. Serre, \emph{Trees}, Springer, 1980.
\bibitem{Vigneras} M.-F. Vign\'eras, \emph{Arithm\'etique des alg\`ebres de quaternions}, Lecture
Notes in Math. 800, Springer, 1980.
\bibitem{Zimmer} R. J. Zimmer, \emph{Ergodic Theory and Semisimple Groups}, Monographs in Math.
81, Birkh\"auser, 1984.
\bibitem{Main} R. Chen, \emph{KMS state uniqueness and phase transitions in localized Bost--Connes
systems: from Chebotarev density to the Hardy--Littlewood conjecture}, preprint, 2026, DOI
\href{https://doi.org/10.5281/zenodo.22152101}{10.5281/zenodo.22152101}.
\end{thebibliography}

\end{document}
